\chapter{Objetivos, hipótesis y factores}
\label{cap:objetivos}

Este capítulo presenta la línea de base, el objetivo superior, el objetivo general técnico, los
objetivos específicos como transformaciones, las hipótesis vinculadas a factores y ratios, y la
cadena visual del procedimiento. Cierra con la figura que enlaza OE1--OE4 con el Capítulo~\ref{cap:metodologia}.

\section{Línea de base / datos iniciales}

En la \textbf{línea de base} de este estudio no se asume un conjunto de mediciones empíricas previas,
sino los artefactos documentales sobre los que opera el modelo de alineación. Son tres piezas que el
procedimiento toma como punto de partida:

\begin{itemize}
  \item \textbf{Entrada principal:} contratos OpenAPI~3.x del sector bancario (objeto de estudio,
        Cap.~1 §1.4).
  \item \textbf{Referencia de contraste:} especificaciones de \ServiceDomain{} BIAN (\Behavior{},
        \BusinessObject{}, relaciones de control).
  \item \textbf{Producto esperado:} representaciones normalizadas, scores S/E/C y
        \AlignmentScore{} con clasificación por niveles.
\end{itemize}

Quedan fuera de esta línea de base las implementaciones en producción, el tráfico \textit{runtime} y
los artefactos de testing automatizado (delimitación §1.6). Los ejemplos del Cap.~1 (Tabla~1.4,
Anexo~F) ilustran el procedimiento pero \textbf{no fijan un caso piloto único}: el diseño aplica a
cualquier par contrato OpenAPI bancario --- \ServiceDomain{} BIAN que cumpla los criterios de
inclusión definidos en el plan.

\section{Objetivo superior}

El \textbf{objetivo superior} expresa la consecuencia esperada en el mundo real \textit{si} el
artefacto existiera y operara en un entorno bancario. No coincide con el objetivo técnico de la
investigación, pero orienta la relevancia del estudio:

\begin{quote}
  Contar con un procedimiento verificable que permita a equipos de arquitectura y gobernanza de APIs
  \textbf{cuantificar la coherencia} entre contratos OpenAPI y modelos BIAN \textbf{antes} del
  despliegue o la certificación, reduciendo retrabajo en adaptadores, ensayo--error en integraciones
  y deriva entre la capa técnica y el modelo de negocio de referencia.
\end{quote}

La investigación \textbf{no implementa} ese cambio en producción (§1.6): entrega el
\textbf{modelo y el diseño del procedimiento} que habilitarían esa consecuencia. La construcción del
artefacto y la validación empírica a escala quedan fuera del alcance del presente plan.

\section{Objetivo general (técnico)}

\textbf{Diseñar un modelo de alineación} que defina cómo calcular \SemanticScore{} (S),
\StructuralScore{} (E), \CoverageScore{} (C) y un \AlignmentScore{} integrado para evaluar contratos
OpenAPI frente a \ServiceDomain{} BIAN, mediante un procedimiento reproducible documentado en este
plan.

\textbf{Base de medición.} El objetivo general se verifica sobre cada instancia
\texttt{contrato OpenAPI bancario $\times$ Service Domain BIAN}: el procedimiento produce
\AlignmentScore{} $\in [0,1]$ y clasificación por niveles (Alta/Media/Baja/Nula). La
\textbf{variable dependiente del artefacto} es el \AlignmentScore{} junto con la clasificación; los
ratios S, E, C y la completitud de normalización son \textbf{productos parciales} de OE1--OE4, no
sustitutos de la VD global.

Cada componente responde a las preguntas P1--P4 del §\ref{sec:planteamiento}: OE1 normaliza los
artefactos; OE2 responde P1 con \textit{matching} estructural (E); OE3 responde P2 con similitud
semántica (S); OE4 integra S, E y C para responder P3 (cobertura) y P4 (score global y
clasificación).

La cadena completa se despliega en la Tabla~\ref{tab:objetivos-especificos}, en la
\Cref{fig:cadena-objetivos} y en la matriz de consistencia (Anexo~C).

\section{Objetivos específicos}

Esta sección formula OE1--OE4 como transformaciones verificables, alineadas con las cinco fases del
modelo de alineación (§\ref{sec:objetivos-metodologia}).

\begin{table}[H]
  \centering
  \caption{Objetivos específicos como transformaciones}
  \label{tab:objetivos-especificos}
  \small
  \begin{tabular}{@{}p{0.08\textwidth}p{0.22\textwidth}p{0.30\textwidth}p{0.30\textwidth}@{}}
    \toprule
    \textbf{ID} & \textbf{Subproblema} & \textbf{Entrada} & \textbf{Salida (producto)} \\
    \midrule
    \textbf{OE1} &
    Adquisición y normalización &
    Contrato OpenAPI bancario + \ServiceDomain{} BIAN &
    Representaciones normalizadas comparables (\textit{paths}, \textit{schemas},
    \Behavior{}, \BusinessObject{}) \\
    \textbf{OE2} &
    Alineación estructural (P1) &
    Representaciones normalizadas &
    Matriz de correspondencias + \StructuralScore{} (E) $\in [0, 1]$ \\
    \textbf{OE3} &
    Alineación semántica (P2) &
    Representaciones normalizadas &
    Matriz de similitud + \SemanticScore{} (S) $\in [0, 1]$ \\
    \textbf{OE4} &
    Evaluación integrada (P3--P4) &
    \StructuralScore{} (E), \SemanticScore{} (S) y representaciones normalizadas (OE1) &
    \CoverageScore{} (C); \AlignmentScore{}; clasificación (Alta/Media/Baja/Nula); tipologías \\
    \bottomrule
  \end{tabular}
\end{table}

La \textbf{Fase~1} del modelo (diseño del procedimiento y definición de scores) corresponde al
marco metodológico general del plan y no se confunde con OE1--OE4, que son
transformadores sucesivos sobre los artefactos.

\section{Hipótesis}

Se plantea una hipótesis por objetivo específico, vinculada al producto verificable y al factor o
ratio que se mediría en una validación posterior (Tabla~\ref{tab:variables}):

\begin{enumerate}
  \item \textbf{H1 (OE1):} Normalizar contratos OpenAPI bancarios y especificaciones BIAN en un
        modelo intermedio común permite calcular de forma sistemática los scores S, E y C
        (\textbf{factor:} representación intermedia; \textbf{ratio:} completitud de la normalización).
  \item \textbf{H2 (OE2):} Las técnicas de \textit{schema matching} (Shvaiko y Euzenat, 2005)
        producen un \StructuralScore{} que refleja la correspondencia entre la organización REST y
        el modelo BIAN de referencia (\textbf{ratio:} E).
  \item \textbf{H3 (OE3):} Las métricas de similitud semántica (El Bayadh et al., 2015;
        Chandrasekaran y Mago, 2022) producen un \SemanticScore{} capaz de distinguir conceptos
        alineados de desalineaciones semánticas o de modelado (\textbf{ratio:} S).
  \item \textbf{H4 (OE4):} La agregación ponderada de S, E y C ($\alpha$, $\beta$, $\gamma$) produce
        un \AlignmentScore{} interpretable y clasificable por niveles de alineación del contrato
        OpenAPI bancario (\textbf{factores:} $\alpha$, $\beta$, $\gamma$; \textbf{ratio:}
        \AlignmentScore{} y clasificación).
\end{enumerate}

Adicionalmente, la \textbf{Fase~1} (diseño) se apoya en la hipótesis de que un modelo estructurado
por componentes S, E y C permite calcular de forma reproducible un score de alineación interpretable
entre OpenAPI y BIAN (detalle en Anexo~C).

\section{Factores y ratios de la investigación}

Los \textbf{factores} son las condiciones de diseño cuya variación puede contrastarse en una
validación empírica posterior; los \textbf{ratios} son los resultados numéricos medidos (scores y
clasificación). Los \textbf{parámetros de control} mantienen constantes otras condiciones para
comparar el efecto de cada factor.

\begin{table}[H]
  \centering
  \caption{Factores y ratios por objetivo específico}
  \label{tab:variables}
  \small
  \begin{tabular}{@{}p{0.08\textwidth}p{0.22\textwidth}p{0.22\textwidth}p{0.36\textwidth}@{}}
    \toprule
    \textbf{OE} & \textbf{Factores} & \textbf{Ratios (resultados)} & \textbf{Parámetros de control} \\
    \midrule
    \textbf{OE1} &
    Esquema de representación intermedia &
    Completitud de extracción (elementos extraídos / esperados) &
    Versión fija del contrato OpenAPI y del \ServiceDomain{} BIAN; herramientas de \textit{parsing} fijas \\
    \textbf{OE2} &
    Estrategia de \textit{schema matching} &
    \StructuralScore{} (E) $\in [0, 1]$ &
    Mismo conjunto normalizado de OE1; umbral de correspondencia fijo \\
    \textbf{OE3} &
    Técnica de similitud semántica (embeddings, ontología, híbrido) &
    \SemanticScore{} (S) $\in [0, 1]$ &
    Mismo corpus de términos; umbral de similitud fijo \\
    \textbf{OE4} &
    Pesos de agregación ($\alpha$, $\beta$, $\gamma$) &
    \CoverageScore{} (C); \AlignmentScore{}; nivel de clasificación &
    $\alpha + \beta + \gamma = 1$ fijos; mismos S, E, C; taxonomía Tabla~\ref{tab:tipologias} \\
    \bottomrule
  \end{tabular}
\end{table}


Los indicadores y métricas detallados por fase se documentan en la matriz de consistencia (Anexo~C).

\section{Figura de cadena de objetivos}

La \Cref{fig:cadena-objetivos} cierra este capítulo: secuencia de transformaciones desde los
artefactos de entrada hasta el \AlignmentScore{} integrado.

\begin{figure}[H]
  \centering
  \includegraphics[width=0.98\textwidth]{Figura_04_Cadena_Objetivos}
  \caption{Cadena de objetivos: normalización $\rightarrow$ scores parciales $\rightarrow$ score global.
           Fase~1 (diseño del modelo) precede a OE1; ver Anexo~C.}
  \label{fig:cadena-objetivos}
\end{figure}
